\section{A remainder for the squared-multiplicity count}
\label{sec:rate}

For positive spacings $a,b$, put $L=a+b$. Use the three nodes
$(0,a,L)$ and weights $(b,-L,a)$. In the normalization of this paper
$\mathcal B_c=a^2 b^2 L/3$. A rate for the variance at all small scales
gives a rate for the sharp count, even for a measure with atoms.

\begin{theorem}[Quantitative inversion for every three-point contrast]
\label{thm:rate}
Let $0<\delta\le1$, $d\ge0$, and suppose that for all sufficiently small $h>0$,
\[
 U_c(h)=d\mathcal B_c h^3\log(1/h)
               +O\bigl(h^3(\log(1/h))^{1-\delta}\bigr).
\]
Then, as $T\to\infty$,
\[
 \frac{\pi A([0,T])}{T\log T}
      =d+O\bigl((\log T)^{-\delta/4}\bigr).
\]
The implied constant can depend on the spacings, $d$, $\delta$, the starting
scale, the bound in the hypothesis, and the mass on a fixed compact interval.
\end{theorem}

\begin{proof}
Write $K=K_c$. Since $K(0)=a^2b^2L^2/4>0$, integration on
$[0,\epsilon_0/h]$ for a fixed small $\epsilon_0$ gives
$A([0,X])\ll (1+X)\log(2+X)$. Compare $A$ with the positive measure
$\mathrm dA_0(x)=(d/\pi)\mathbf1_{[1,\infty)}(x)\log x\,\mathrm dx$.
The estimates $K(x)\ll\min(1,x^{-4})$ give, by $y=hx$,
\[
 U_{c,A_0}(h)=d\mathcal B_c h^3\log(1/h)+O(h^3).
\]
Indeed $\int_0^\infty K(y)|\log y|\,\mathrm dy<\infty$ and the contribution of
$0<y<h$ after rescaling is $O(h^4)$, since
$\int_0^h|\log(y/h)|\,\mathrm dy=h$.
For the signed difference $\sigma=A-A_0$, put
\[
 R(t)=\int_0^\infty
       \frac{\cos(tx)-1+(tx)^2/2}{x^4}\,\mathrm d\sigma(x),
 \qquad r(t)=R(t)/t^3\quad(t>0).
\]
The continuous value of the integrand at zero is $t^4/24$.
The numerator is $O(\min(|tx|^4,|tx|^2))$. Splitting at $1/t$ and
integrating the tail by parts using the count bound gives
$|r(t)|\ll1+\log(1/t)$ at small $t$. For $t\ge t_0>0$,
$|R(t)|\ll_{t_0}t^4$, by splitting instead at $x=1$.

Expanding the variance cancels the constant and quadratic terms. If
$u(h)=(U_c(h)-U_{c,A_0}(h))/h^3$, then
\begin{equation}\label{eq:rate-contract}
 r(t)=\frac{u(t/L)}{2abL^3}
    +\frac{a^2}{L^2}r(at/L)+\frac{b^2}{L^2}r(bt/L).
\end{equation}
The two recursion coefficients have sum $\rho=(a^2+b^2)/L^2<1$.
Put $p=1-\delta$ and $C_0=\log(L/\min(a,b))$.
After $j$ substitutions the new forcing terms are bounded in total by
a constant times $\rho^j(1+\log(1/t)+j C_0)^p$, for $t$ below a fixed
sufficiently small $t_0$. The unreplaced remainder at depth $j$ is
$O(\rho^j(1+\log(1/t)+j C_0))$ and tends to zero. Summation proves
\[
 |R(t)|\ll t^3(1+\log(1/t))^{1-\delta}\qquad(0<t<t_0).
\]

To pass to an interval, set
$\mathrm d\sigma_T(x)=\pi\,\mathrm d\sigma(Tx)/(T\log T)$ in pushforward
notation. Its corresponding transform is
$R_T(t)=\pi T^3R(t/T)/\log T$. For $1\le |t|\le t_0T$,
\[
 |R_T(t)|\ll |t|^3(\log T)^{-\delta}.
\]
For $|t|<1$, the same upper bound times
$1+|\log|t||$ suffices, by evenness and the estimate on $R$. For
$|t|\ge t_0T$, use $|R_T(t)|\ll |t|^4/(T\log T)$.

Let $f$ be an even smooth compactly supported real function and use
$\widehat g(t)=\int_{\mathbb R}g(x)e^{-itx}\,\mathrm dx$. Fourier inversion gives
\begin{equation}\label{eq:rate-invert}
 \int_0^\infty f(x)\,\mathrm d\sigma_T(x)
 =\frac1{2\pi}\int_{\mathbb R} R_T(t)\,\widehat{x^4 f}(t)\,\mathrm dt .
\end{equation}
In fact, $g(x)=x^4 f(x)$ has $g(0)=g''(0)=0$, so subtracting the constant
and quadratic Taylor terms from the cosine leaves its inverse Fourier
integral unchanged. At zero the formula follows by the fourth derivative.
Fubini follows by the same numerator bound applied to the total variation
of $\sigma_T$ and the rapid decay of $\widehat g$.

Choose a fixed even nonnegative smooth unit-mass mollifier supported in
$[-1,1]$, scaled to width $\epsilon<1/4$. Convolve it with the indicators of
$[-1+\epsilon,1-\epsilon]$ and $[-1-\epsilon,1+\epsilon]$ to obtain $f_-,f_+$
sandwiching $\mathbf1_{[-1,1]}$. For either choice and any integer $N\ge1$,
\[
 |\widehat{x^4 f_\pm}(t)|\ll_N (1+|t|)^{-1}(1+\epsilon |t|)^{-N}.
\]
The first derivative of $x^4f_\pm$ has bounded $L^1$ norm and its derivative
of order $N+1$ has $L^1$ norm $O_N(\epsilon^{-N})$. Both estimates follow by
differentiating the mollified endpoints and integrating by parts.

In \eqref{eq:rate-invert} the integral for $|t|\le t_0T$ is consequently
$O((\log T)^{-\delta}\epsilon^{-3})$, from the integral of
$t^2(1+\epsilon t)^{-N}$ over $t\ge1$ with $N>3$. For
$\epsilon=(\log T)^{-\delta/4}$, $N=8$ makes the remaining tail
$O(\epsilon^{-8}T^{-5}/\log T)$.
The two smooth tests against the rescaling of $A_0$ differ from its
interval count by $O(\epsilon)$, since its normalized density is bounded
on $[1-2\epsilon,1+2\epsilon]$. Positivity of $A$ sandwiches its interval
count between the two tests. Its difference from the normalized $A_0$ count
is $O(\epsilon+(\log T)^{-\delta}\epsilon^{-3})$.
Finally $\pi A_0([0,T])/(T\log T)=d+O(1/\log T)$, which gives the assertion.
\end{proof}

Under the repository's RH input the prime-error variance has
$A=2\sum_{\gamma>0}m_\gamma^2\delta_\gamma$ and the denominator
$(1/4+\gamma^2)^2$ in place of $\gamma^4$. The substitution costs $O(h^4)$:
$|p_c(hx)|\ll h^2 x^2$ and
$1-x^4/(1/4+x^2)^2\ll(1+x^2)^{-1}$. The latter is integrable against
$A$ by the classical unit-interval zero bound $O(\log(2+x))$ counting
multiplicities, which gives $O(\log^2(2+x))$ for their squares. Thus a measured
remainder $O(h^3)$ gives
\[
 2\pi\sum_{0<\gamma\le T}^{\rm distinct}m_\gamma^2
   =dT\log T+O\bigl(T(\log T)^{3/4}\bigr).
\]
This implication provides a count error for any supplied variance remainder.
Establishing that remainder for the actual zeta zeros is a further arithmetic
problem. The exponent $\delta/4$ comes from smoothing an interval of relative
width $\epsilon$ and is not claimed optimal. For background on Fourier
Tauberian estimates for monotone counting functions, see
\href{https://doi.org/10.1006/jfan.2001.3764}{Yu.~Safarov, J.~Funct.\ Anal.\ 185
(2001), 111--128}.
